% ==========================================================================
%  Section 4 -- The physical-volume covariance limit.  Owner: W04.
%  Frozen label: sec:covariance.  Target 1200 words.
%
%  Contents: the two pairings (eq:pairings); the periodic massive resolvent
%  built from an energy form and a Riesz adjoint (lem:resolvent); the
%  physical-cell Riemann-sum theorem (ext:boxintegral); the lattice
%  covariance as a counting pairing with the inverse of the Gaussian operator
%  (lem:resolvent-pairing); the fixed-volume limit in five labelled steps
%  (thm:covariance-limit); and the absence of a Bessel kernel (rem:no-bessel).
%
%  Naming: the pairings are written Qcal_N, Qcal, because the letter Q is the
%  Gaussian operator m^2 - Delta_a of def:lattice-laplacian (section 3).
% ==========================================================================
\section{The physical-volume covariance limit}\label{sec:covariance}

This section proves the one analytic limit of Part~I that is not measure theory:
the covariance integrals of the physical lattice laws converge, and their limit
is a pairing with a resolvent constructed before any Fourier series is written.
\Cref{lem:resolvent} builds that resolvent from an energy form and a Riesz
representation \cite{yosida1980}, \cref{lem:resolvent-pairing} rewrites the
lattice covariance of
\cref{def:gj-measure} as a counting pairing with the inverse of the Gaussian
operator $A_N=m^2-\Delta_a$, and \cref{thm:covariance-limit} combines the two.
The objects are the pairings
\begin{equation}\label{eq:pairings}
  \mathcal Q_N(f,g):=\frac1V\sum_{p\in I_N}
    \frac{\overline{F_N(p)}\,G_N(p)}{\lambda_N(p)},
  \qquad
  \mathcal Q(f,g):=\frac1V\sum_{k\in\Z^2}
    \frac{\overline{F(k)}\,G(k)}{\lambda(k)},
\end{equation}
where $F$ is \eqref{eq:fourier-coefficient}, $F_N$ the sampled coefficient
\eqref{eq:sampled-coefficient}, and $\lambda_N,\lambda$ the symbols
\eqref{eq:lattice-symbol}--\eqref{eq:continuum-symbol}. The calligraphic letter
keeps $Q$ for the Gaussian operator of \cref{def:lattice-laplacian}, and no
other letter is free. Since
$\overline{F_N(p)}=F_N(-p)$ and both symbols are even, the sums are real for
real $f,g$; the computation is dimension-blind except in one shell count,
flagged in Step~5 below, and we write the two-dimensional case that
\cref{sec:setting} fixes.

\begin{lemma}[The periodic massive resolvent]\label{lem:resolvent}
Let $\nu$ be Lebesgue measure on $[0,1)^2$, of total mass $1$, and
$E:=L^2([0,1)^2,\nu;\C)$ with $\ip{u}{v}_E=\int\overline{u}\,v\,\dd\nu$,
conjugate-linear in the first argument. Write $b_k(z)=e^{2\pi i\,k\cdot z}$ for
the basis of \cref{ext:fourier-basis} and $B:E\to H$ for the induced isometry
onto $H:=\ell^2(\Z^2;\C)$. Put $\lambda(k)=m^2+\sum_i(2\pi k_i/L)^2$ and
$w(k):=\lambda(k)^{-1/2}$, and let $J$ and the gradient multipliers $D_i$ act
diagonally by $(Jh)(k)=w(k)h(k)$ and
$(D_ih)(k)=i\,(2\pi k_i/L)\,w(k)h(k)$; the value and gradient maps are
$jh:=B^{-1}Jh$ and $B^{-1}D_ih$. Then:
\begin{enumerate}
\item[(i)] the sesquilinear form
  \begin{equation}\label{eq:energy-form}
    \mathcal E(h,v):=m^2\ip{jh}{jv}_E+\sum_{i=1}^2\ip{D_ih}{D_iv}_E
  \end{equation}
  is the energy form of $H^1(\Torus)$ and satisfies
  $\mathcal E(h,v)=\ip{h}{v}_H$ for all $h,v\in H$, and the physical-volume
  inner product on the torus of period $L$ is $V\ip{\cdot}{\cdot}_E$;
\item[(ii)] $\{(jh,(D_ih)_i):h\in H\}$ is the closure in $E\times E^2$ of the
  pairs formed by finite trigonometric polynomials on $\Torus$ and their
  classical physical-coordinate derivatives; this is the periodic
  $H^1(\Torus)$ graph completion of the energy form in (i);
\item[(iii)] for $g\in E$ the vector $h_g:=J^*Bg$, $Rg:=jh_g$, is the unique
  element of $E$ satisfying
  \begin{equation}\label{eq:weak-equation}
    m^2\ip{Rg}{jv}_E+\sum_{i=1}^2\ip{D_ih_g}{D_iv}_E=\ip{g}{jv}_E
  \end{equation}
  for every $v\in H$, and $\norm{Rg}_E\le m^{-2}\norm{g}_E$;
\item[(iv)] $B(Rg)(k)=\lambda(k)^{-1}Bg(k)$, and
  $\ip{f}{Rg}_E=\sum_k\overline{Bf(k)}\,Bg(k)/\lambda(k)$, the series being
  absolutely convergent.
\end{enumerate}
The construction is carried out on the unit-mass $L^2$ of the torus. No
distribution-defined Sobolev space is invoked: the completion in (ii) is an
explicit graph closure of polynomials.
\end{lemma}

\begin{proof}
(i) The multipliers are diagonal in the basis $b_k$, and
\[
  m^2w(k)^2+\sum_i(2\pi k_i/L)^2w(k)^2=\lambda(k)w(k)^2=1
\]
gives
\eqref{eq:energy-form}. The identity concerns the mixed sesquilinear form, not
only its diagonal; physical volume enters only through $\dd x=V\dd\nu$.
(ii) $J$ is bounded and injective because $0<w(k)\le m^{-1}$, and
finite-support vectors are dense in $H$. A finite-support $h$ gives the
polynomial $\sum_ka(k)b_k$ with $a(k)=w(k)h(k)$, whose coordinate derivative is
$\sum_ki\,(2\pi k_i/L)\,a(k)b_k$; conversely every finite-support polynomial
arises this way, as $w(k)\ne0$. If value-gradient pairs are Cauchy in
$E\times E^2$, \eqref{eq:energy-form} makes the energy vectors Cauchy,
completeness of $H$ supplies their limit, and continuity of $j$ and $D_i$
identifies the graph limit; the range is closed.
(iii) Since $\norm{jv}_E\le m^{-1}\norm{v}_H$, the functional
$v\mapsto\ip{g}{jv}_E$ is bounded on $H$ and represented by $h_g=J^*Bg$, and
$\ip{h_g}{v}_H=\ip{g}{jv}_E$ with \eqref{eq:energy-form} gives
\eqref{eq:weak-equation}. Two solutions differ by $u$ with
$\mathcal E(u,v)=0$ for every $v$; testing $v=u$ gives $u=0$, and injectivity of
$j$ fixes the $E$-value. Testing \eqref{eq:weak-equation} with $v=h_g$ and
taking real parts gives
$m^2\norm{Rg}_E^2\le\Real\ip{g}{Rg}_E\le\norm{g}_E\norm{Rg}_E$, and we divide
when $\norm{Rg}_E\ne0$.
(iv) The adjoint identity on a single basis vector gives $(J^*c)(k)=w(k)c(k)$
because $w$ is real, so $B(Rg)(k)=w(k)^2Bg(k)=\lambda(k)^{-1}Bg(k)$;
Cauchy--Schwarz and $\lambda\ge m^2$ make the coefficient series absolutely
convergent.
\end{proof}

\begin{external}[Physical-cell Riemann sums converge]\label{ext:boxintegral}
Let $g:[0,L]^2\to\C$ be continuous, with sup-norm modulus of continuity
$\omega_g$, and partition $[0,L]^2$ into the $N^2$ cubes of side $a=L/N$. The
sum of the lower-corner sample times $a^2$ over all cubes differs from
$\int_{[0,L]^2}g\,\dd x$ by at most $V\omega_g(a)$, so these Riemann sums
converge; half-open cells give the same value, the shared boundaries being
Lebesgue null. This is the Henstock--Kurzweil box-integral theory of the
standard library, quoted and not re-proved.
\end{external}

\begin{lemma}[Lattice covariance versus resolvent pairing]\label{lem:resolvent-pairing}
Let $b,c\in\R^{I_N}$ and let $f,g$ be real $L$-periodic. Then
\begin{equation}\label{eq:counting-pairing}
  \ip{C_Nb}{C_Nc}=a^{-2}\ip{b}{A_N^{-1}c}_{\mathrm{counting}},
  \qquad
  a^2\ip{f}{A_N^{-1}g}_{\mathrm{counting}}=\mathcal Q_N(f,g).
\end{equation}
\end{lemma}

\begin{proof}
The first identity is \eqref{eq:gj-pairing} of \cref{def:gj-measure}, the
inverse-mass form in which the lattice covariance is used: $C_N$ is an input
object (\cref{rem:lattice-input}), and the power $a^{-2}$ sits in the pairing,
not in a rescaling of the field. For the second identity, use the
character eigenbasis of \cref{prop:operator}: a real lattice function $u$ has
$u(x)=V^{-1}\sum_pF_N(p)\chi_p(x)$ because $N^2a^2=L^2=V$, and
$Q\chi_p=\lambda_N(p)\chi_p$ is \eqref{eq:eigen-identity}, with $A_N=Q$ acting
in the counting inner product. Substituting both expansions into the
bilinear counting pairing and using the orthogonality
$\sum_{x\in I_N}\chi_p(x)\chi_q(x)=N^2\mathbf 1_{p+q\equiv0}$ together with
$G_N(-p)=\overline{G_N(p)}$ for real $g$ gives
\[
  a^2\ip{f}{A_N^{-1}g}_{\mathrm{counting}}
  =a^2V^{-2}N^2\sum_pF_N(p)\,\overline{G_N(p)}\,\lambda_N(p)^{-1}
  =\overline{\mathcal Q_N(f,g)}=\mathcal Q_N(f,g),
\]
the last equality because $\mathcal Q_N$ is real for real $f,g$. The powers of
$a$ are not cancelled by hand: each sampled vector carries $a^2$ against the
$a^{-2}$ of the pairing, and $N^2a^2=V$ turns the counting sum into the
normalised pairing. Together with \eqref{eq:sampled-covariance}, this writes the
covariance pairing in the symbol form the limit needs.
\end{proof}

The measure below is the physical lattice law $\mu^{\mathrm{GJ}}_N$ of
\cref{def:gj-measure}, not the finite-source pushforward $\mu_N$ whose laws
converge in \cref{sec:tightness}--\cref{sec:main}: what is proved here is a
statement about pairings, and it enters \cref{sec:marginals} as the pairwise
covariance input, since every finite cylinder coordinate of the lift is a finite
linear combination of lattice source coordinates (\cref{def:finite-source-lift});
the entrywise identification is separate (\cref{lem:hC-identification}).

\begin{theorem}[Fixed-volume covariance limit]\label{thm:covariance-limit}
Let $L>0$ and $m>0$, and let $f,g:\R^2\to\R$ be bounded continuous and
$L$-periodic. Let $a=L/N$, let $\mu^{\mathrm{GJ}}_N$ be the physical lattice law
of \cref{def:gj-measure}, write $f_L\in E$ for $z\mapsto f(Lz)$ as in
\cref{lem:haar-lebesgue}, and let $b,c\in\R^{I_N}$ be the sampled vectors
$b_x=a^2f(ax)$, $c_x=a^2g(ax)$. Then, as $N\to\infty$ with $L$ and $m$ fixed,
\begin{equation}\label{eq:three-limits}
  \begin{aligned}
    \ip{C_Nb}{C_Nc}&\longrightarrow\Real V\ip{f_L}{Rg_L}_E,\\
    \int\omega(f)\,\omega(g)\,\mu^{\mathrm{GJ}}_N(\dd\omega)
      &\longrightarrow\Real V\ip{f_L}{Rg_L}_E,\\
    \int\bigl(\omega(f)-\E\omega(f)\bigr)
         \bigl(\omega(g)-\E\omega(g)\bigr)\,\mu^{\mathrm{GJ}}_N(\dd\omega)
      &\longrightarrow\Real V\ip{f_L}{Rg_L}_E,
  \end{aligned}
\end{equation}
where $\omega(f)$ is the test average of \cref{def:gj-measure} and $\E$ is the
expectation under $\mu^{\mathrm{GJ}}_N$. Both means vanish identically in $N$, so
the second and third limits coincide.
\end{theorem}

\begin{proof}
By \cref{lem:resolvent-pairing} the first line of \eqref{eq:three-limits} is
$\mathcal Q_N(f,g)$, and the other two are the mixed second moment and the
centred covariance of the same Gaussian family. It suffices to prove
$\mathcal Q_N(f,g)\to\mathcal Q(f,g)$ and to identify $\mathcal Q$ with the
resolvent pairing, in five labelled steps.

\textbf{Step 1 (fixed-mode Riemann sums).} Let $k\in\Z^2$ and let
$p_N(k)\in I_N$ be its residue class. At the sampled physical point $t=a x$ the
phase is exactly $e^{-2\pi i\,k\cdot t/L}=e^{-2\pi i\,k\cdot x/N}$, for either
representative of the residue class, and
$t\mapsto f(t)e^{-2\pi i\,k\cdot t/L}$ is uniformly continuous on the compact
cell with sup-norm modulus of continuity $\omega_k$. \Cref{ext:boxintegral}
applied to that function gives
\begin{equation}\label{eq:fixed-mode}
  \abs{F_N(p_N(k))-F(k)}\le V\omega_k(a),
  \qquad \omega_k(a)\longrightarrow0.
\end{equation}
Continuity on the compact cell is what produces $\omega_k(a)\to0$; boundedness
alone would not, since a bounded function may oscillate across every cube of
every mesh.

\textbf{Step 2 (dispersion).} The limit $\sin t/t\to1$ gives
$4(N/L)^2\sin^2(\pi k_i/N)\to(2\pi k_i/L)^2$, hence
$\lambda_N(p_N(k))\to\lambda(k)$. Both symbols are at least $m^2>0$, so their
products and quotients converge; and the residue map is injective on
$\{k:\norm{k}_\infty\le K\}$ once $N>2K$, so for fixed $K$ the sum over the
cutoff converges term by term, real parts included.

\textbf{Step 3 (uniform denominators and tails).} Choose for each residue its
centered representative in
$\{-\lfloor N/2\rfloor,\dots,\lceil N/2\rceil-1\}^2$, which represents the
even-$N$ Nyquist frequency once. The chord bound $\sin t\ge2t/\pi$ on
$[0,\pi/2]$ gives, for centered $k$,
\begin{equation}\label{eq:uniform-denominator}
  \lambda_N(k)\ge m^2+\frac{16\norm{k}_\infty^2}{L^2},
  \qquad\text{so that}\qquad
  R_K:=m^2+\frac{16K^2}{L^2}
\end{equation}
bounds $\lambda_N$ once $\norm{k}_\infty>K$. Write $\mathcal Q_{N,K}$ and
$\mathcal Q_K$ for the sums \eqref{eq:pairings} restricted to
$\norm{k}_\infty\le K$. Cauchy--Schwarz with the finite Parseval identity
(\cref{lem:finite-parseval}) --- where boundedness of $f$ enters, and the only
place where it does --- gives
\begin{equation}\label{eq:discrete-tail}
  \abs{\mathcal Q_N-\mathcal Q_{N,K}}\le\frac{1}{VR_K}
    \Bigl(\sum_p\abs{F_N(p)}^2\Bigr)^{1/2}
    \Bigl(\sum_p\abs{G_N(p)}^2\Bigr)^{1/2}
  \le\frac{V\norm{f}_\infty\norm{g}_\infty}{R_K}
\end{equation}
uniformly in $N$ and $K\ge0$; when no modes remain outside the cutoff the left
side is zero. On the continuum side $\norm{k}_\infty\ge K+1$ gives
$\lambda(k)\ge D_K:=m^2+(2\pi/L)^2(K+1)^2$, and the same Cauchy--Schwarz
estimate on the annuli between symmetric boxes gives
\begin{equation}\label{eq:continuum-tail}
  \abs{\mathcal Q-\mathcal Q_K}
  \le\frac{V\norm{f}_\infty\norm{g}_\infty}{D_K}\longrightarrow0 .
\end{equation}
Because $4\pi^2>16$ the constant $D_K$ is at least $R_K$, so
\eqref{eq:continuum-tail} yields the tail with the constant of the discrete
estimate; that inequality is what lets one cutoff serve both sides.

\textbf{Step 4 (eventual Bessel bound).} Fix the box
$S_K=\{k:\norm{k}_\infty\le K\}$. For $N>2K$ reduction modulo $N$ is injective
on $S_K$ and identifies $S_K$ with the discrete cutoff, so nonnegativity and
\cref{lem:finite-parseval} give
\begin{equation}\label{eq:bessel-eventual}
  \sum_{k\in S_K}\abs{F_N(p_N(k))}^2\le V^2\norm{f}_\infty^2
  \qquad (N>2K).
\end{equation}
The restriction $N>2K$ is not a convenience. The bound is \emph{false} as a
statement about all $N$: a small grid aliases distinct integer modes onto one
residue, so the left side counts one coefficient several times and can exceed
the right side without bound. Since every finite $S\subseteq\Z^2$ lies in some
$S_K$, the fixed-mode limits \eqref{eq:fixed-mode} in
\eqref{eq:bessel-eventual} give
$\sum_{k\in S}\abs{F(k)}^2\le V^2\norm{f}_\infty^2$ for every finite $S$, and
bounded finite sums of a nonnegative series give square summability of
$k\mapsto F(k)$. No continuous Bessel theorem and no Fourier completeness is
assumed: this is the Bessel \emph{inequality}, and it is finite-dimensional.

\textbf{Step 5 (closed form, summability, envelope, assembly).} By (iii)--(iv)
of \cref{lem:resolvent} and the Haar-to-Lebesgue normalisation
$Bf_L(k)=V^{-1}F(k)$ of \cref{lem:haar-lebesgue},
\begin{equation}\label{eq:closed-form}
  V\ip{f_L}{Rg_L}_E
  =V\sum_k\overline{V^{-1}F(k)}\;\frac{V^{-1}G(k)}{\lambda(k)}
  =V^{-1}\sum_k\frac{\overline{F(k)}\,G(k)}{\lambda(k)}
  =\mathcal Q(f,g).
\end{equation}
So the Fourier Green pairing already carries the physical factor $L^{-2}$, and
the identification of $\mathcal Q$ with the resolvent pairing is a computation,
not a definition of $\mathcal Q$ as an operator. Absolute summability follows
from Young's inequality and $\lambda\ge m^2$,
\begin{equation}\label{eq:young}
  \abs{\frac{\overline{F(k)}\,G(k)}{\lambda(k)}}
  \le\frac{\abs{F(k)}^2+\abs{G(k)}^2}{m^2},
\end{equation}
so the series in \eqref{eq:pairings} converges, and exhaustion by symmetric
boxes gives $\mathcal Q_K\to\mathcal Q$.

A modewise majorant independent of $N$ needs more than continuity. Fix $r=d+1$,
which is $r=3$ for the case $d=2$ treated here, and
$D_r(f):=\max_i\norm{\partial_i^r f}_\infty$. The $r$-fold forward difference at
spacing $a$ is bounded by $a^rD_r(f)$ and multiplies the sampled sum by
$(\chi-1)^r$; on a nonzero centered frequency with $\norm{k}_\infty=q$ the chord
bound gives $\abs{e^{2\pi i k_i/N}-1}\ge4q/N$, and
$1/\lambda_N,1/\lambda\le L^2/(16q^2)$ yields the common envelope
\begin{equation}\label{eq:envelope}
  \abs{F_N(k)}\le V\Bigl(\frac L4\Bigr)^{r}D_r(f)\,q^{-r},
  \qquad
  \frac1V\abs{\frac{\overline{F_N(k)}\,G_N(k)}{\lambda_N(k)}}
  \le\frac{VL^2}{16}\Bigl(\frac L4\Bigr)^{2r}D_r(f)D_r(g)\,q^{-2r-2}.
\end{equation}
The sup-norm shell $\norm{k}_\infty=q$ carries $(2q+1)^d-(2q-1)^d$ modes, at
most $2d(2q+1)^{d-1}$; in $d=2$ the count is $8q$ (\cref{ext:summability}).
Since $2r+2>d$, the shell majorant is summable, and the zero mode is bounded by
$m^{-2}$ and $V\norm{f}_\infty\norm{g}_\infty$. The envelope \eqref{eq:envelope}
is a statement about $C^{r}$ periodic tests. It is \emph{not} available for all
continuous tests, and no such statement is claimed: the limit below uses no
envelope, which is why \cref{thm:covariance-limit} concerns bounded continuous
periodic tests and is not a modewise convergence theorem.

Choose $K$ so that both tails \eqref{eq:discrete-tail} and
\eqref{eq:continuum-tail} are below $\epsilon/3$; Step~2 then bounds the
fixed-cutoff error by $\epsilon/3$ for all large $N$. The triangle inequality
gives
\begin{equation}\label{eq:epsilon-three}
  \abs{\Real\mathcal Q_N(f,g)-\Real\mathcal Q(f,g)}
  \le\abs{\Real\mathcal Q_N-\Real\mathcal Q_{N,K}}
   +\abs{\Real\mathcal Q_{N,K}-\Real\mathcal Q_K}
   +\abs{\Real\mathcal Q_K-\Real\mathcal Q}
  <\epsilon,
\end{equation}
with one $K$ serving all $N$, and with no caller supplying summability, a tail
estimate or a convergence hypothesis. Transferring \eqref{eq:epsilon-three} to
the measures uses only that $\mu^{\mathrm{GJ}}_N$ is Gaussian with covariance
$C_N$ (\cref{def:gj-measure}) and has mean zero exactly at every mesh: then the
mixed second moment and the centred covariance integral both equal
$\Real\mathcal Q_N$ and hence converge to $\Real V\ip{f_L}{Rg_L}_E$ by
\eqref{eq:closed-form}. This proves \eqref{eq:three-limits}.
\end{proof}

\begin{remark}[Why no Bessel kernel appears]\label{rem:no-bessel}
The proof of \cref{thm:covariance-limit} introduces no lattice Green function
$G_a(x-y)$, no Bessel function, no modified Bessel function $K_0$, no constant
$1/(2\pi)$ or $1/(4\pi)$, and no logarithmic divergence. Every estimate is an
estimate on symbols: $\lambda_N^{-1}$ on $I_N$ and $\lambda^{-1}$ on $\Z^2$,
compared in \eqref{eq:uniform-denominator}--\eqref{eq:continuum-tail} and
controlled by \eqref{eq:bessel-eventual} and \eqref{eq:young}. Nothing is lost
by staying in the symbol domain: what is proved is a convergence statement
about pairings of bounded continuous periodic tests with the resolvent, and such
a pairing is determined by the spectral data of the resolvent on the Fourier
basis, which is exactly $\lambda^{-1}$; a positional kernel is that same data
written in another coordinate system, and producing it in closed form is a
special-function computation the argument neither performs nor needs. A text
that wrote the two-dimensional kernel $C_m(x)=K_0(m\abs{x})/(2\pi)$ here would be
naming the $x$-space representative of the \emph{same} operator on the same torus ---
the two are equal by the classical Green-function identity, and evaluating
\eqref{eq:pairings} against that kernel reproduces $\Real Q$ --- but that
identity is neither used above nor proved above, and this paper does not prove
it; what is proved is the limit of the pairings, in the symbol domain, with the
constants of this section.
Nor is a logarithmic divergence at issue: at fixed volume
$\lambda_N(0)=\lambda(0)=m^2$, and the high-frequency behaviour is absorbed by
\eqref{eq:young} through the test coefficients rather than by a kernel.
\end{remark}

With \cref{thm:covariance-limit} the analysis of the lattice symbol is finished.
Every later step of Part~I is measure theory on the carrier of
\cref{sec:carrier}: the uniform moment and compact tail of
\cref{sec:tightness}, the finite marginals of \cref{sec:marginals}, and the
assembly in \cref{sec:main}. No further symbol estimate occurs, and none is
needed.
